% GENERATED FILE. Edit canonical lessons, then run npm run generate:books.
% canonical-source-hash: fnv1a64:51065458c35b48ec
% canonical-lessons: FR-C502-flou, FR-C502-net, FR-C502-cliche, FR-C502-cadrer, FR-C502-trepied

\chapter{Prendre des photos}
\label{ch:fr-prendre-des-photos}
\hlchaptermodality{502}

\begin{chapteropening}
\textbf{By the end of this chapter:} \emph{I can talk about taking good photos.}

Complete the last lesson of chapter 502: I can talk about taking good photos.
\end{chapteropening}

\section[flou]{flou — blurred, out of focus; vague}

\label{lesson:FR-C502-flou}



\begin{quote}
Before the new one: say the French for model-making, then the French for to put together.
\end{quote}

\subsection*{You'll want to know: flou}
\textbf{flou} — \textquotedblleft{}blurred, out of focus; vague\textquotedblright{}.

\begin{quote}
Toutes mes photos de la fête sont floues.
\end{quote}

\begin{quote}
\emph{All my photos from the party are blurred.}
\end{quote}

Feminine \emph{floue}. Also vague: \emph{des idées floues}.

\begin{cousinweb}
\emph{flou} goes back to Old French \emph{flou} \textquotedblleft{}faded, weak, soft\textquotedblright{} (EtymDB). Its earlier origin is not settled.
\end{cousinweb}

\begin{grammarlens}[title={What you've built}]
One more word for telling, explaining and describing at B1.
\end{grammarlens}

\subsection*{Your turn}
\begin{itemize}
\raggedright
  \item \emph{Say it:} \emph{flou}
  \item \emph{Say it:} \emph{flou}, once more
  \item \emph{Say it:} a sentence of your own with \emph{flou}
  \item \emph{Recall:} say the French for model-making, then the French for to put together, then say \emph{flou} again
\end{itemize}

\begin{tcolorbox}[breakable,colback=teal!4,colframe=teal!35!black,title={Before you move on}]
What does flou mean? (\textquotedblleft{}Blurred, out of focus; vague\textquotedblright{}.) Say it once more.
\end{tcolorbox}
\section[net]{net — sharp, in focus; clean, clear}

\label{lesson:FR-C502-net}



\begin{quote}
Before the new one: say the French for to put together, then the French for blurred.
\end{quote}

\subsection*{You'll want to know: net}
\textbf{net} — \textquotedblleft{}sharp, in focus; clean, clear\textquotedblright{}.

\begin{quote}
Sur cette photo, le visage est bien net.
\end{quote}

\begin{quote}
\emph{In this photo the face is nice and sharp.}
\end{quote}

Feminine \emph{nette}. Also net (of money): \emph{un salaire net}.

\begin{cousinweb}
\emph{net} is inherited from Latin \emph{nitidus} \textquotedblleft{}shining, clean\textquotedblright{} (EtymDB), from \emph{nitēre} \textquotedblleft{}to shine\textquotedblright{} (Lewis \& Short).
\end{cousinweb}

\begin{grammarlens}[title={What you've built}]
One more word for telling, explaining and describing at B1.
\end{grammarlens}

\subsection*{Your turn}
\begin{itemize}
\raggedright
  \item \emph{Say it:} \emph{net}
  \item \emph{Say it:} \emph{net}, once more
  \item \emph{Say it:} a sentence of your own with \emph{net}
  \item \emph{Recall:} say the French for to put together, then the French for blurred, then say \emph{net} again
\end{itemize}

\begin{tcolorbox}[breakable,colback=teal!4,colframe=teal!35!black,title={Before you move on}]
What does net mean? (\textquotedblleft{}Sharp, in focus; clean, clear\textquotedblright{}.) Say it once more.
\end{tcolorbox}
\section[le cliché]{le cliché — snapshot, photo; cliché}

\label{lesson:FR-C502-cliche}



\begin{quote}
Before the new one: say the French for blurred, then the French for sharp.
\end{quote}

\subsection*{You'll want to know: le cliché}
\textbf{le cliché} — \textquotedblleft{}snapshot, photo; cliché\textquotedblright{}.

\begin{quote}
Le photographe a pris de beaux clichés du port.
\end{quote}

\begin{quote}
\emph{The photographer took some lovely shots of the harbour.}
\end{quote}

Also a worn-out idea, as in English.

\begin{cousinweb}
\emph{le cliché} is the participle of \emph{clicher} \textquotedblleft{}to make a printing plate\textquotedblright{} (EtymDB), a word probably imitating the sound of the die striking the metal.
\end{cousinweb}

\begin{grammarlens}[title={What you've built}]
One more word for telling, explaining and describing at B1.
\end{grammarlens}

\subsection*{Your turn}
\begin{itemize}
\raggedright
  \item \emph{Say it:} \emph{le cliché}
  \item \emph{Say it:} \emph{le cliché}, once more
  \item \emph{Say it:} a sentence of your own with \emph{le cliché}
  \item \emph{Recall:} say the French for blurred, then the French for sharp, then say \emph{le cliché} again
\end{itemize}

\begin{tcolorbox}[breakable,colback=teal!4,colframe=teal!35!black,title={Before you move on}]
What does le cliché mean? (\textquotedblleft{}Snapshot, photo; cliché\textquotedblright{}.) Say it once more.
\end{tcolorbox}
\section[cadrer]{cadrer — to frame (a picture)}

\label{lesson:FR-C502-cadrer}



\begin{quote}
Before the new one: say the French for sharp, then the French for snapshot.
\end{quote}

\subsection*{You'll want to know: cadrer}
\textbf{cadrer} — \textquotedblleft{}to frame (a picture)\textquotedblright{}.

\begin{quote}
Elle a mal cadré la photo : la tour est coupée.
\end{quote}

\begin{quote}
\emph{She framed the photo badly: the tower is cut off.}
\end{quote}

Also to fit: \emph{ça ne cadre pas avec ses idées}.

\begin{cousinweb}
\emph{cadrer} is built on \emph{cadre} \textquotedblleft{}frame\textquotedblright{}, borrowed from Italian \emph{quadro}, from Latin \emph{quadrum} \textquotedblleft{}a square\textquotedblright{} (Lewis \& Short).
\end{cousinweb}

\begin{grammarlens}[title={What you've built}]
One more word for telling, explaining and describing at B1.
\end{grammarlens}

\subsection*{Your turn}
\begin{itemize}
\raggedright
  \item \emph{Say it:} \emph{cadrer}
  \item \emph{Say it:} \emph{cadrer}, once more
  \item \emph{Say it:} a sentence of your own with \emph{cadrer}
  \item \emph{Recall:} say the French for sharp, then the French for snapshot, then say \emph{cadrer} again
\end{itemize}

\begin{tcolorbox}[breakable,colback=teal!4,colframe=teal!35!black,title={Before you move on}]
What does cadrer mean? (\textquotedblleft{}To frame (a picture)\textquotedblright{}.) Say it once more.
\end{tcolorbox}
\section[le trépied]{le trépied — tripod}

\label{lesson:FR-C502-trepied}



\begin{quote}
Before the new one: say the French for snapshot, then the French for to frame.
\end{quote}

\subsection*{You'll want to know: le trépied}
\textbf{le trépied} — \textquotedblleft{}tripod\textquotedblright{}.

\begin{quote}
Pour les photos de nuit, il utilise un trépied.
\end{quote}

\begin{quote}
\emph{For night photos he uses a tripod.}
\end{quote}

Three feet: \emph{tré-} \textquotedblleft{}three\textquotedblright{} and \emph{pied} \textquotedblleft{}foot\textquotedblright{}.

\begin{cousinweb}
\emph{le trépied} comes from Latin \emph{tripēs} \textquotedblleft{}three-footed\textquotedblright{} (EtymDB), built from \emph{trēs} \textquotedblleft{}three\textquotedblright{} and \emph{pēs} \textquotedblleft{}foot\textquotedblright{} (Lewis \& Short).
\end{cousinweb}

\begin{grammarlens}[title={What you've built}]
One more word for telling, explaining and describing at B1.
\end{grammarlens}

\subsection*{Your turn}
\begin{itemize}
\raggedright
  \item \emph{Say it:} \emph{le trépied}
  \item \emph{Say it:} \emph{le trépied}, once more
  \item \emph{Say it:} a sentence of your own with \emph{le trépied}
  \item \emph{Recall:} say the French for snapshot, then the French for to frame, then say \emph{le trépied} again
\end{itemize}

\begin{tcolorbox}[breakable,colback=teal!4,colframe=teal!35!black,title={Before you move on}]
What does le trépied mean? (\textquotedblleft{}Tripod\textquotedblright{}.) Say it once more.
\end{tcolorbox}
